% Đề xuất Project 3 --- biên dịch được bằng pdfLaTeX (mặc định Overleaf) hoặc XeLaTeX/LuaLaTeX.
% Các chỗ cần điền đánh dấu bằng \dien{...} (hiện màu đỏ khi biên dịch).
\documentclass[11pt,a4paper]{article}

\usepackage{iftex}
\ifPDFTeX
  \usepackage[T5]{fontenc}
  \usepackage[utf8]{inputenc}
  \usepackage{lmodern}
\else
  \usepackage{fontspec}
\fi
\usepackage[vietnamese]{babel}

\usepackage[a4paper,margin=2.2cm]{geometry}
\usepackage{microtype}
\usepackage{xcolor}
\usepackage{booktabs}
\usepackage{tabularx}
\usepackage{array}
\usepackage{enumitem}
\usepackage{tikz}
\usetikzlibrary{arrows.meta,positioning,fit,calc}
\usepackage{caption}
\usepackage{fancyhdr}
\usepackage[hidelinks]{hyperref}

\definecolor{accent}{HTML}{0F766E}
\definecolor{ink}{HTML}{1F2937}
\definecolor{muted}{HTML}{6B7280}
\definecolor{boxbg}{HTML}{F3F4F6}
\definecolor{newbg}{HTML}{ECFDF5}
\definecolor{oldbg}{HTML}{F3F4F6}
\definecolor{extbg}{HTML}{FEF3C7}

\setlist{itemsep=2pt,topsep=4pt}
\setlength{\parskip}{4pt}
\setlength{\parindent}{0pt}
\renewcommand{\arraystretch}{1.25}
\newcolumntype{Y}{>{\raggedright\arraybackslash}X}

\pagestyle{fancy}
\fancyhf{}
\lhead{\small\color{muted} Project 3 --- Đề xuất đề tài}
\rhead{\small\color{muted} \thepage}
\renewcommand{\headrulewidth}{0.4pt}

\newcommand{\code}[1]{\texttt{#1}}
\newcommand{\dien}[1]{\textcolor{red}{[#1]}}
\newcommand{\keybox}[2]{%
  \par\noindent\colorbox{boxbg}{\parbox{\dimexpr\linewidth-2\fboxsep}{\vspace{3pt}\textbf{#1}\par\vspace{2pt}#2\vspace{3pt}}}\par}

\begin{document}

\begin{center}
  {\small\color{muted} ĐẠI HỌC BÁCH KHOA HÀ NỘI --- TRƯỜNG CÔNG NGHỆ THÔNG TIN VÀ TRUYỀN THÔNG}\\[2pt]
  {\small\color{muted} KHOA KHOA HỌC MÁY TÍNH}\\[10pt]
  {\Large\bfseries\color{ink} ĐỀ XUẤT ĐỀ TÀI PROJECT 3}\\[8pt]
  {\large\bfseries Agent tự bồi đắp kho tri thức cho chatbot nội bộ trên nền DeepSeek Harness:\\
  xử lý câu hỏi ngoài khả năng của RAG, tìm nguồn ngoài có kiểm duyệt\\ và giải thích lựa chọn}\\[4pt]
  {\normalsize\itshape A Self-Enriching Knowledge Agent for Enterprise Chatbots on DeepSeek Harness}
\end{center}
\vspace{6pt}

\begin{tabularx}{\linewidth}{@{}lY@{}}
  \toprule
  Sinh viên thực hiện & \dien{Họ và tên} --- MSSV \dien{...} --- Lớp \dien{...} \\
  Giảng viên hướng dẫn & \dien{Học hàm, học vị, họ và tên} \\
  Thời gian & 12 tuần, từ \dien{ngày bắt đầu} đến \dien{ngày kết thúc}; báo cáo tiến độ 2 tuần một lần (6 mốc) \\
  Hướng nghiên cứu & Agentic AI, agentic RAG, thiết kế harness cho LLM agent, con người trong vòng lặp \\
  \bottomrule
\end{tabularx}

\vspace{6pt}
\keybox{Tóm tắt}{%
Chatbot nội bộ dạng RAG chỉ trả lời được những gì có sẵn trong kho tài liệu. Câu hỏi nào trượt sẽ thành
phiếu hỗ trợ (ticket) và quản trị viên phải xử lý tay. Đề tài xây dựng một \textbf{agent} trên DeepSeek
Harness (DSH) nhận chính các ticket đó: \textbf{chẩn đoán} vì sao chatbot trượt, \textbf{tự tìm lại} trong
kho nội bộ (kể cả tính toán bằng mã trong sandbox), \textbf{tìm nguồn ngoài} qua một agent con bị giới hạn
khi kho thật sự thiếu, rồi soạn \textbf{phiếu đề xuất} kèm \textbf{lý do chọn nguồn} để quản trị viên duyệt.
Chỉ sau khi được duyệt, câu trả lời mới gửi cho người hỏi và kho tri thức mới được cập nhật, nhờ đó
chatbot tự trả lời được các câu tương tự về sau.
\begin{itemize}[leftmargin=14pt]
  \item \textbf{Sản phẩm:} agent xử lý ticket (plugin và preset DSH), agent con tìm nguồn ngoài, giao diện
        quản trị viên duyệt đề xuất và cập nhật kho.
  \item \textbf{Nghiên cứu:} bộ câu hỏi tiếng Việt có đáp án và nhãn nguyên nhân trượt, web giả lập có nguồn
        xấu, so sánh 5 cấu hình harness H0--H4, và thí nghiệm đánh giá tác dụng của phiếu lý do đối với
        người duyệt.
  \item \textbf{Nền tảng kế thừa:} hệ thống đa người dùng có sandbox mà sinh viên đã xây dựng trước
        (mục~\ref{sec:nen-tang}). Toàn bộ 12 tuần dành cho agent, dữ liệu và đánh giá.
\end{itemize}}

% ================================================================ 1
\section{Đặt vấn đề}

\subsection{Chatbot RAG nội bộ và giới hạn của nó}
Sinh viên đã xây dựng DocShield, một chatbot hỏi đáp tài liệu nội bộ theo kiểu RAG~\cite{rag} có chống ảo
giác: tìm kiếm kết hợp vector (bge-m3~\cite{bgem3}) và BM25, bắt buộc trích dẫn, mã nguồn (không phải mô
hình) kiểm tra trích dẫn khớp nguyên văn trước khi hiển thị, và từ chối kèm tạo ticket khi không tìm thấy.
Cách làm này trả lời tốt câu hỏi có sẵn trong một đoạn tài liệu, nhưng có ba điểm yếu:
\begin{itemize}
  \item \textbf{Nhiều ticket thực ra trả lời được.} Câu hỏi cần tổng hợp nhiều tài liệu, cần tính toán, hoặc
        dùng từ khác với tài liệu, nên tìm kiếm một lượt bị trượt.
  \item \textbf{Quản trị viên tốn thời gian} cho các việc lặp lại: tự tra, tự trả lời, tự bổ sung tài liệu.
  \item \textbf{Chatbot không học được gì:} câu hỏi đó lần sau vẫn trượt.
\end{itemize}

\subsection{Khi nào cần agent}
RAG một lượt là đủ khi câu trả lời nằm sẵn trong một đoạn văn. Agent (mô hình gọi công cụ trong vòng lặp
suy luận và hành động~\cite{react}) chỉ đáng dùng khi câu hỏi cần nhiều bước, ví dụ:
\begin{itemize}
  \item \textbf{Tổng hợp nhiều tài liệu:} ``Muốn chuyển ngành thì cần điều kiện gì và mất bao nhiêu học phí?''
  \item \textbf{Tính toán trên bảng:} ``Với điểm này em có đủ học bổng loại Giỏi không?''
  \item \textbf{Đối chiếu phiên bản:} ``Quy chế mới khác quy chế cũ ở điểm nào?''
  \item \textbf{Kho thiếu thông tin:} quy định mới của cơ quan cấp trên chưa được đưa vào kho.
\end{itemize}
Vì agent đắt và khó kiểm soát hơn, đề tài dùng \textbf{kiến trúc lai}~\cite{anthropic-agents}: luồng RAG
cố định (workflow) vẫn trả lời câu dễ, nhanh và rẻ; agent chỉ vào cuộc khi workflow trượt.

\subsection{Vấn đề mở}
Cho agent tự do hơn, đặc biệt cho phép lấy dữ liệu từ bên ngoài, đặt ra ba câu hỏi:
\begin{enumerate}
  \item Agent trả lời được thêm bao nhiêu câu, và \textbf{có làm tăng tỉ lệ bịa} không? Chống ảo giác vốn
        là mục tiêu của hệ thống, nên trở thành ràng buộc phải giữ.
  \item Những \textbf{cơ chế của harness} (chạy mã, lập kế hoạch, agent con, workflow lập trình được) đóng
        góp bao nhiêu, và tốn thêm bao nhiêu?
  \item Nguồn ngoài có thể kém tin cậy, đã cũ, hoặc chứa lệnh ẩn (prompt injection gián
        tiếp~\cite{greshake}). Làm sao để \textbf{người duyệt bắt được đề xuất sai} mà không phải tự
        kiểm tra lại từ đầu?
\end{enumerate}

% ================================================================ 2
\section{Mục tiêu và câu hỏi nghiên cứu}

\subsection{Mục tiêu sản phẩm}
\begin{enumerate}[label=P\arabic*.]
  \item Agent xử lý ticket: chẩn đoán nguyên nhân trượt, tự tìm lại trong kho nội bộ, tính toán bằng mã
        trong sandbox.
  \item Agent con tìm nguồn ngoài có giới hạn: có internet nhưng không thấy dữ liệu nội bộ, chỉ truy cập tên
        miền được phép, lưu bản chụp từng nguồn.
  \item Phiếu đề xuất cho quản trị viên: câu trả lời có trích dẫn kiểm chứng, đề xuất cập nhật kho, và
        \textbf{phiếu lý do chọn nguồn}.
  \item Giao diện duyệt: duyệt, sửa hoặc từ chối; khi duyệt thì trả lời người hỏi và cập nhật kho.
\end{enumerate}

\subsection{Câu hỏi nghiên cứu}
\begin{description}[leftmargin=14pt,style=nextline]
  \item[RQ1 --- Hiệu quả.] Kiến trúc lai trả lời đúng được thêm bao nhiêu câu so với RAG thuần, tốn thêm bao
        nhiêu token và thời gian, và tỉ lệ bịa có tăng không?
  \item[RQ2 --- Thiết kế harness.] Mỗi cơ chế của DSH (chạy mã, plan mode, subagent, workflow) đóng góp bao
        nhiêu vào độ chính xác? Loại câu hỏi nào nên bật cơ chế nào?
  \item[RQ3 --- Con người trong vòng lặp.] Phiếu lý do chọn nguồn có giúp người duyệt phát hiện đề xuất sai
        (nguồn kém, số liệu cũ, trang có lệnh ẩn) \textbf{đúng hơn và nhanh hơn} không?
\end{description}

% ================================================================ 3
\section{Nền tảng kế thừa và phần việc mới}\label{sec:nen-tang}

Đề tài không xây lại hạ tầng. Các thành phần dưới đây đã được sinh viên xây dựng và kiểm thử trước khi bắt
đầu Project 3 (134 test tự động đều đạt); chúng chỉ được dùng lại.

\begin{tabularx}{\linewidth}{@{}>{\raggedright\arraybackslash}p{0.47\linewidth}Y@{}}
\toprule
\textbf{Đã có (kế thừa)} & \textbf{Làm mới trong Project 3} \\
\midrule
\begin{itemize}[leftmargin=12pt,nosep]
  \item dsh-gate: đăng nhập, phiên, quản trị, mỗi người dùng một DSH
  \item DocShield: RAG lai, trích dẫn được mã kiểm tra, luồng ticket
  \item Sandbox S0--S4: container riêng, DB tách ra dịch vụ, proxy giữ API key, tường lửa, Landlock
  \item Agent có shell, đọc/ghi file, subagent trong sandbox
\end{itemize}
&
\begin{itemize}[leftmargin=12pt,nosep]
  \item Agent xử lý ticket và các công cụ chẩn đoán
  \item Agent con tìm nguồn ngoài và proxy web có kiểm soát
  \item Phiếu đề xuất, phiếu lý do, giao diện duyệt
  \item Bộ dữ liệu, web giả lập, bộ chấm tự động
  \item Thí nghiệm H0--H4 và thí nghiệm người duyệt
\end{itemize} \\
\bottomrule
\end{tabularx}

% ================================================================ 4
\section{Phương pháp}

\subsection{Kiến trúc lai}

\begin{center}
\begin{tikzpicture}[
  box/.style={draw, rounded corners=3pt, align=center, font=\small, inner sep=5pt, text width=3.1cm},
  arr/.style={-{Stealth[length=2.2mm]}, thick, color=muted},
  lbl/.style={font=\scriptsize, color=muted}]
  \node[box, fill=boxbg] (u) {Người dùng hỏi};
  \node[box, fill=oldbg, right=0.9cm of u] (w) {\textbf{Workflow RAG}\\(đã có)};
  \node[box, fill=boxbg, right=0.9cm of w] (ok) {Trả lời có\\trích dẫn};
  \node[box, fill=newbg, below=0.9cm of w] (a) {\textbf{Agent xử lý ticket}\\chẩn đoán, tìm lại, chạy mã};
  \node[box, fill=oldbg, left=0.9cm of a] (kb) {Kho nội bộ\\(DocShield)};
  \node[box, fill=extbg, right=0.9cm of a] (ext) {\textbf{Agent con}\\nguồn ngoài có giới hạn};
  \node[box, fill=newbg, below=0.9cm of a] (p) {\textbf{Phiếu đề xuất}\\+ lý do chọn nguồn};
  \node[box, fill=boxbg, below=0.9cm of p] (ad) {\textbf{Quản trị viên duyệt}};
  \node[box, fill=boxbg, left=0.9cm of ad] (ans) {Trả lời người hỏi};
  \node[box, fill=boxbg, right=0.9cm of ad] (up) {Cập nhật kho};
  \draw[arr] (u) -- (w);
  \draw[arr] (w) -- node[above, lbl]{tìm thấy} (ok);
  \draw[arr] (w) -- node[right, lbl]{trượt $\rightarrow$ ticket} (a);
  \draw[arr, <->] (a) -- (kb);
  \draw[arr, <->] (a) -- node[above, lbl]{khi kho thiếu} (ext);
  \draw[arr] (a) -- (p);
  \draw[arr] (p) -- (ad);
  \draw[arr] (ad) -- (ans);
  \draw[arr] (ad) -- (up);
  \draw[arr, dashed] (up.north) |- ($(kb.south)+(0,-0.35)$) -- (kb.south);
\end{tikzpicture}
\end{center}

\subsection{Quy trình của agent}
\begin{enumerate}
  \item \textbf{Chẩn đoán} vì sao workflow trượt, gán một trong các nhãn: tìm kiếm trượt (tài liệu có nhưng
        dùng từ khác), cần nhiều bước (tổng hợp, tính toán), tài liệu thiếu, tài liệu cũ hoặc mâu thuẫn,
        câu hỏi mơ hồ.
  \item \textbf{Tự tìm lại trong kho nội bộ:} viết lại câu hỏi, tra nhiều tài liệu, đọc bảng và tính toán
        bằng mã trong sandbox. Mọi trích dẫn vẫn phải qua bộ kiểm tra khớp nguyên văn của DocShield.
  \item \textbf{Tìm nguồn ngoài khi kho thiếu,} bằng cách giao việc cho agent con (mục~\ref{sec:an-toan}).
  \item \textbf{Soạn phiếu đề xuất} gồm: câu trả lời có trích dẫn, hoặc kết luận ``tài liệu thật sự không
        có''; đề xuất cập nhật kho (thêm tài liệu, đánh dấu tài liệu cũ, thêm cách hỏi khác cho đoạn đã có);
        và phiếu lý do chọn nguồn.
  \item \textbf{Quản trị viên duyệt, sửa hoặc từ chối.} Không có gì được gửi cho người dùng hay ghi vào kho
        trước bước này.
\end{enumerate}

\subsection{Phiếu lý do chọn nguồn}
Với mỗi nguồn ngoài được đề xuất, agent phải điền một phiếu có cấu trúc cố định:
\begin{itemize}
  \item \textbf{Lấp chỗ trống nào:} câu hỏi hoặc phần kho đang thiếu mà nguồn này bổ sung.
  \item \textbf{Độ tin cậy:} loại nguồn (văn bản chính thức, trang tổ chức, diễn đàn\ldots) và lý do.
  \item \textbf{Độ mới:} ngày ban hành hoặc cập nhật, và có văn bản nào thay thế không.
  \item \textbf{Các nguồn đã xem nhưng loại,} kèm lý do loại.
  \item \textbf{Xung đột với kho hiện có:} đoạn nào trong kho nói khác, và vì sao đề xuất ưu tiên nguồn mới.
\end{itemize}
Phiếu giúp người duyệt kiểm tra lập luận của agent thay vì phải tự tìm kiếm lại. RQ3 đo xem điều này có
đúng không.

\subsection{An toàn khi lấy dữ liệu ngoài}\label{sec:an-toan}
Tận dụng sandbox sẵn có:
\begin{itemize}
  \item \textbf{Tách quyền:} agent con có internet nhưng chạy trong container riêng, \textbf{không thấy kho nội
        bộ và dữ liệu người dùng}; agent chính thấy kho nhưng không có internet. Nhờ vậy một trang web độc hại
        không thể khiến agent làm lộ dữ liệu nội bộ.
  \item \textbf{Đi qua proxy web:} chỉ truy cập tên miền trong danh sách được phép; mọi trang được lưu bản chụp
        để người duyệt xem lại đúng nội dung agent đã đọc.
  \item \textbf{Nội dung ngoài là dữ liệu, không phải lệnh:} agent con chỉ trả về trích đoạn kèm địa chỉ; web
        giả lập có cài lệnh ẩn để đo tỉ lệ agent bị lừa.
\end{itemize}

% ================================================================ 5
\section{Thiết kế thí nghiệm}

\subsection{Dữ liệu}
\begin{itemize}
  \item \textbf{Kho tài liệu:} tài liệu nội bộ tiếng Việt dạng quy chế, quy định, biểu mẫu \dien{nguồn cụ thể}.
  \item \textbf{Bộ câu hỏi:} 150--200 câu, mỗi câu có đáp án, đoạn bằng chứng và nhãn ``vì sao trượt''. Một phần
        lấy từ log ticket thật của DocShield, phần còn lại tự xây dựng để đủ các nhãn. Chia tập dev và tập test;
        tập test chỉ chạy một lần ở cuối.
  \item \textbf{Web giả lập:} gồm trang chính thức, trang kém tin cậy, trang đã cũ và trang có lệnh ẩn, để kết
        quả lặp lại được và đo được khả năng bị nguồn xấu đánh lừa.
\end{itemize}

\subsection{Các cấu hình so sánh}
\begin{tabularx}{\linewidth}{@{}lY@{}}
\toprule
\textbf{Cấu hình} & \textbf{Mô tả} \\
\midrule
H0 & RAG thuần (DocShield hiện tại): mốc so sánh \\
H1 & Agent + tìm lại nhiều lượt + chạy mã trong sandbox \\
H2 & H1 + plan mode (lập kế hoạch trước khi làm) \\
H3 & H2 + subagent (chia việc cho agent con trong kho nội bộ) \\
H4 & H3 + agent con tìm nguồn ngoài có kiểm duyệt \\
\bottomrule
\end{tabularx}

Mỗi câu chạy 3--5 lần cho mỗi cấu hình, báo cáo trung bình và độ lệch để giảm ảnh hưởng của tính ngẫu
nhiên. Thí nghiệm RQ3: người duyệt \dien{số người} xem cùng một tập phiếu đề xuất (có trộn đề xuất sai), một
nửa có phiếu lý do và một nửa không, đổi chéo giữa người duyệt.

\subsection{Chỉ số}
\begin{tabularx}{\linewidth}{@{}>{\raggedright\arraybackslash}p{0.34\linewidth}Y@{}}
\toprule
\textbf{Chỉ số} & \textbf{Cách đo} \\
\midrule
Tỉ lệ trả lời đúng (RQ1, RQ2) & So với đáp án; chấm tự động kết hợp LLM làm giám khảo~\cite{llmjudge}, kiểm tra tay một mẫu \\
Tỉ lệ bịa (RQ1) & Tỉ lệ câu có trích dẫn không kiểm chứng được hoặc khẳng định không có trong nguồn \\
Độ chính xác chẩn đoán & Nhãn agent gán so với nhãn ``vì sao trượt'' \\
Tỉ lệ bị nguồn xấu lừa & Tỉ lệ đề xuất dùng trang kém tin cậy, đã cũ hoặc làm theo lệnh ẩn \\
Hiệu quả bồi đắp & Sau khi cập nhật kho, tỉ lệ câu tương tự mà H0 tự trả lời được \\
Người duyệt (RQ3) & Tỉ lệ phát hiện đề xuất sai, thời gian ra quyết định \\
Chi phí & Token, số lượt gọi công cụ, thời gian mỗi ticket \\
\bottomrule
\end{tabularx}

% ================================================================ 6
\section{Phạm vi, giới hạn và rủi ro}

\textbf{Trong phạm vi:} tài liệu nội bộ tiếng Việt; bộ câu hỏi 150--200 câu; cấu hình H0--H4; web giả lập; luồng
quản trị viên duyệt.

\textbf{Ngoài phạm vi:} agent tự gửi câu trả lời hoặc tự sửa kho mà không qua duyệt; truy cập internet không
giới hạn; dữ liệu có cấu trúc (cơ sở dữ liệu, bảng số liệu lớn).

\begin{tabularx}{\linewidth}{@{}>{\raggedright\arraybackslash}p{0.32\linewidth}Y@{}}
\toprule
\textbf{Rủi ro} & \textbf{Cách giảm} \\
\midrule
Bộ câu hỏi nhỏ, thiên lệch & Cân bằng theo nhãn; báo cáo kết quả theo từng nhãn; tập test tách riêng \\
Ít người duyệt cho RQ3 & Thiết kế đổi chéo; bổ sung phân tích định tính \\
Chi phí token cao khi chạy lặp & Chạy tập dev nhỏ trước; dùng mô hình rẻ cho bộ chấm \\
Web giả lập khác web thật & Dựng từ bản chụp trang thật; thử thêm một phần nhỏ trên web thật có danh sách tên miền \\
DSH còn là bản alpha & Cố định phiên bản trong image; ghi lại phiên bản trong mọi kết quả \\
\bottomrule
\end{tabularx}

% ================================================================ 7
\section{Kế hoạch thực hiện và các mốc báo cáo}

\begin{tabularx}{\linewidth}{@{}c c Y >{\raggedright\arraybackslash}p{0.27\linewidth}@{}}
\toprule
\textbf{Mốc} & \textbf{Tuần} & \textbf{Nội dung} & \textbf{Demo / minh chứng} \\
\midrule
1 & 1--2 & Tổng quan (agentic RAG~\cite{agentic-rag-survey}, Self-RAG~\cite{selfrag}, CRAG~\cite{crag},
ReAct, Reflexion~\cite{reflexion}, prompt injection); chốt RQ; kiến trúc; thu thập tài liệu &
Kiến trúc + phân loại các kiểu trượt từ log thật của DocShield \\
2 & 3--4 & Bộ câu hỏi có đáp án, web giả lập, bộ chấm tự động; chạy H0 & Bảng số liệu mốc so sánh \\
3 & 5--6 & Agent v1: chẩn đoán, tìm lại, chạy mã; chạy H1--H3 trên tập dev & Demo agent xử lý một ticket khó \\
4 & 7--8 & Agent con tìm nguồn ngoài (proxy, danh sách tên miền, bản chụp); phiếu lý do; chạy H4 &
Demo phiếu đề xuất có nguồn ngoài \\
5 & 9--10 & Giao diện duyệt, cập nhật kho, chạy lại câu cũ; thí nghiệm RQ3; chạy tập test & Bảng kết quả chính \\
6 & 11--12 & Phân tích lỗi và chi phí; báo cáo, slide, video & Nộp \\
\bottomrule
\end{tabularx}

% ================================================================ 8
\section{Sản phẩm bàn giao}
\begin{enumerate}
  \item Agent xử lý ticket và agent con tìm nguồn ngoài (plugin và preset DSH), tích hợp với DocShield và gate.
  \item Giao diện quản trị viên duyệt đề xuất và cập nhật kho.
  \item Bộ câu hỏi tiếng Việt có nhãn, web giả lập, bộ đánh giá tự động.
  \item Báo cáo kết quả H0--H4 và RQ3, kèm khuyến nghị cấu hình harness cho từng loại câu hỏi.
  \item Mã nguồn, hướng dẫn chạy lại thí nghiệm, slide và video demo.
\end{enumerate}

% ================================================================ 9
\section{Hướng phát triển thành đồ án tốt nghiệp}
Project 3 trả lời câu hỏi ``agent làm được tới đâu và cần người duyệt ra sao''. Đồ án tốt nghiệp dự kiến đi
tiếp theo hướng \textbf{agent học dần từ quyết định của quản trị viên}, nối thẳng từ RQ3:
\begin{itemize}
  \item Lưu các lần duyệt và từ chối làm \textbf{bộ nhớ kinh nghiệm}, dùng để học \textbf{chính sách chọn
        nguồn} (nguồn nào hay bị từ chối, vì sao).
  \item \textbf{Giảm dần việc của người duyệt:} đề xuất đủ tin cậy thì tự áp dụng có kiểm soát và có thể hoàn
        tác; đo khi nào tự động hóa là an toàn.
  \item \textbf{Giữ kho tri thức ``sống'':} tự phát hiện tài liệu hết hạn và mâu thuẫn.
\end{itemize}
Nếu kết quả Project 3 thiên về sản phẩm hơn, phương án thay thế là mở rộng sang \textbf{nhiều loại nguồn}:
thêm dữ liệu có cấu trúc qua lớp ngữ nghĩa (semantic layer) và nhiều agent chuyên trách có agent điều phối.

% ================================================================ TLTK
\begin{thebibliography}{99}\small
\bibitem{rag} P.~Lewis et al. \emph{Retrieval-Augmented Generation for Knowledge-Intensive NLP Tasks}. NeurIPS, 2020.
\bibitem{react} S.~Yao, J.~Zhao, D.~Yu, N.~Du, I.~Shafran, K.~Narasimhan, Y.~Cao. \emph{ReAct: Synergizing Reasoning and Acting in Language Models}. ICLR, 2023.
\bibitem{reflexion} N.~Shinn, F.~Cassano, A.~Gopinath, K.~Narasimhan, S.~Yao. \emph{Reflexion: Language Agents with Verbal Reinforcement Learning}. NeurIPS, 2023.
\bibitem{selfrag} A.~Asai, Z.~Wu, Y.~Wang, A.~Sil, H.~Hajishirzi. \emph{Self-RAG: Learning to Retrieve, Generate, and Critique through Self-Reflection}. ICLR, 2024.
\bibitem{crag} S.-Q.~Yan, J.-C.~Gu, Y.~Zhu, Z.-H.~Ling. \emph{Corrective Retrieval Augmented Generation}. arXiv:2401.15884, 2024.
\bibitem{agentic-rag-survey} A.~Singh, A.~Ehtesham, S.~Kumar, T.~T.~Khoei. \emph{Agentic Retrieval-Augmented Generation: A Survey on Agentic RAG}. arXiv:2501.09136, 2025.
\bibitem{anthropic-agents} Anthropic. \emph{Building Effective Agents}. Bài viết kỹ thuật, 12/2024.
\bibitem{greshake} K.~Greshake et al. \emph{Not What You've Signed Up For: Compromising Real-World LLM-Integrated Applications with Indirect Prompt Injection}. AISec Workshop, 2023.
\bibitem{bgem3} J.~Chen et al. \emph{BGE M3-Embedding: Multi-Lingual, Multi-Functionality, Multi-Granularity Text Embeddings Through Self-Knowledge Distillation}. arXiv:2402.03216, 2024.
\bibitem{llmjudge} L.~Zheng et al. \emph{Judging LLM-as-a-Judge with MT-Bench and Chatbot Arena}. NeurIPS Datasets and Benchmarks, 2023.
\end{thebibliography}

\end{document}
